% % Appendix Chapter

% Appendix 1 \hskip 2.5em Example Abstract 
% \vspace{7em}

% \noindent Note 1: Page numbering is done by consecutive numbering. Additional appendices, that are attached to the thesis (such as copies, drawings, etc.) are left without page numbers and placed at the end of the thesis.

% \newpage
% \markboth{left}{\normalfont{Appendix 1. Example Abstract}}

% \noindent \textbf{Jurmu M.\ (2007) Resource Management in Smart Spaces Using Context-Based Leases.} University of Oulu, Department of Electrical and Information Engineering. Master’s Thesis, 77 p.

% \begin{center}
% \textbf{\fontsize{16}{19pt}\selectfont ABSTRACT}\\
% \end{center}

% {\bfseries \noindent The convergence of wireless access networks in conjunction with the increased computing power of the handheld terminals is preparing the emergence of the ubiquitous computing paradigm. In the center of this paradigm are smart spaces, which are local environments saturated with various embedded computational resources. These spaces co-operate with mobile client devices in enabling advanced, service-oriented computation scenarios. This co-operation is typically enabled through the utilization of distributed and modular middleware frameworks. An emerging additional requirement however is the possibility to harness resources from the proximity environment to the mobile device in an on-demand fashion. This is a challenge especially to the resource management infrastructure of the smart spaces.

% This thesis explores the concept of a smart space and presents a review of the current research and technologies. Subsequently, a lease-based design for resource management in smart spaces is presented. Leases in this work are negotiated agreements between the mobile clients and the resource management infrastructure, regarding the harnessed resources. Leasing is seen as a suitable solution for the management due to the transient nature of the resource usage. The inclusion of additional contextual features to the leases further facilitates the management.

% Requirements for the design are derived from the review and from an example usage scenario. The requirements include dynamic mapping and contracting of resources from the proximity environment, monitoring of the contract validity, access control towards the resources and dynamic maintenance of the smart space infrastructure. Presented design is analytically compared against existing solutions, and several points for future development are listed. According to the comparison, none of the existing solutions utilize contracts with contextual validity in resource management. Two publications of this work have been accepted into an international conference and a workshop on the focus area of pervasive computing.

% \vspace{1\baselineskip}
% \noindent\textbf{Keywords: Ubiquitous computing, mobile computing, task-based computing, context-awareness, QoS.}
% }

% Auto-generated from converted markdown.
% Source: converted_markdown/appendices.md

% \# Chapter 7 Appendices

% This appendix chapter consolidates supplementary material that supports the methodological and reproducibility claims in the main thesis. The appendices are intended to document materials that are important for transparency but too detailed to place in the core chapter narrative.

% \section{Appendix Contents}

% The appendices should contain, at minimum, the following materials:

% 1. prompt templates used in the main thesis-facing comparisons; 1. model identifiers, providers, and decoding settings where available; 1. benchmark and comparison artifact paths; 1. annotation instructions and tone-label definitions; 1. additional difficult examples and disagreement cases; 1. supplementary failure-mode tables; 1. extended reproducibility notes for rerunning stored analyses.

This appendix chapter consolidates supplementary material that supports the methodological and reproducibility claims in the main thesis. The appendices are intended to document materials that are important for transparency but too detailed to place in the core chapter narrative.

\section{Appendix Contents}\label{appendix-contents}

The appendices should contain, at minimum, the following materials:

\begin{enumerate}
\def\labelenumi{\arabic{enumi}.}
\tightlist
\item
  prompt templates used in the main thesis-facing comparisons;
\item
  model identifiers, providers, and decoding settings where available;
\item
  benchmark and comparison artifact paths;
\item
  annotation instructions and tone-label definitions;
\item
  additional difficult examples and disagreement cases;
\item
  supplementary failure-mode tables;
\item
  extended reproducibility notes for rerunning stored analyses.
\end{enumerate}

\section{Tone-Label Definitions}\label{tone-label-definitions}

The pilot annotation and review process should use the following working definitions:

\begin{enumerate}
\def\labelenumi{\arabic{enumi}.}
\tightlist
\item
  \textbf{Commitment:} forward-looking promises, targets, intentions, or policy declarations that do not yet demonstrate completed implementation or realized outcome.
\item
  \textbf{Action:} concrete activities, agreements, programs, investments, or operational steps that are being undertaken or have been initiated.
\item
  \textbf{Outcome:} realized, completed, or demonstrably achieved results, milestones, or measured performance changes.
\item
  \textbf{Needs review:} records whose language is too ambiguous, mixed, or incomplete for stable tone assignment without additional human interpretation.
\end{enumerate}

\section{Reproducibility Notes}\label{reproducibility-notes}

The main reproducibility artifacts referenced in the thesis include:

\begin{enumerate}
\def\labelenumi{\arabic{enumi}.}
\tightlist
\item
  OCR-expanded document folders under \texttt{data/thesis\_dataset/};
\item
  extracted ESG records under \texttt{results/esg\_records.json};
\item
  resumable benchmark outputs under \texttt{results/t1\_results.jsonl} and \texttt{results/t2\_results.jsonl};
\item
  prompt and model stability summaries under \texttt{results/revision\_analysis/};
\item
  dashboard-ready figures and tables under \texttt{results/thesis\_workflow\_dashboard/}.
\end{enumerate}

These artifacts support rerunning, auditing, and extending the workflow even where exact third-party LLM outputs may vary across time.

\section{End-to-End User Workflow Diagram}\label{end-to-end-user-workflow-diagram}

The appendix should also document the operational user flow across the main thesis-facing tools. In the implemented repository, the workflow begins with OCR preparation in \texttt{pages/Bulk\_OCR.py}, continues through page-range-aware ESG extraction in \texttt{pages/llm\_processing.py}, and then passes through the ClimateBERT comparison stage as the T1 pipeline.

\subsection{Workflow Summary}\label{workflow-summary}

The implemented user flow is:

\begin{enumerate}
\def\labelenumi{\arabic{enumi}.}
\tightlist
\item
  upload or select PDF files in \textbf{Bulk OCR}
\item
  run OCR and save page-level outputs into \texttt{data/thesis\_dataset/}
\item
  open \textbf{LLM Processing}
\item
  select the OCR-expanded PDF/document target
\item
  choose a page range for processing
\item
  choose one of three LLM-provider families:

  \begin{itemize}
  \tightlist
  \item
    \textbf{OpenRouter}
  \item
    \textbf{LM Studio / OpenAI-compatible}
  \item
    \textbf{Ollama}
  \end{itemize}
\item
  run structured ESG extraction over the selected page range
\item
  save structured records into \texttt{results/esg\_records.json}
\item
  run \textbf{ClimateBERT} processing as the T1 comparison layer
\item
  save ClimateBERT outputs into \texttt{results/predictions.json} or resumable T1 artifacts
\end{enumerate}


\begin{figure}[H]
\begin{center}
  \includegraphics*[width=\textwidth]{mermaid_flow_diagram_01}
\end{center}
\caption{Bulk OCR to LLM-processing flow diagram. This figure visualizes the operational path from file intake, OCR expansion, and page storage to provider-specific ESG extraction and downstream ClimateBERT comparison outputs.}
\alt{Flowchart showing user entry through Bulk OCR, OCR output storage, document selection in LLM Processing, provider choice among OpenRouter, LM Studio, and Ollama, ESG record storage, and downstream ClimateBERT output generation.}
\label{fig:commitment_outcome_ratio}
\end{figure}

\subsection{Mermaid Sequence Diagram}\label{mermaid-sequence-diagram}

\begin{figure}[H]
\begin{center}
  \includegraphics*[width=\textwidth]{mermaid_flow_diagram_02}
\end{center}
\caption{Bulk OCR to ClimateBERT sequence diagram. This figure shows the interaction order between the user, OCR storage, LLM processing, provider backends, ESG record storage, and the downstream ClimateBERT comparison stage.}
\alt{Sequence diagram showing the user uploading or selecting files, OCR artifact creation, page-range selection in LLM Processing, requests to an LLM provider, ESG record persistence, and ClimateBERT prediction storage.}
\label{fig:appendix_bulkocr_llm_climatebert_sequence}
\end{figure} 

Sequence explanation:

\begin{enumerate}
\def\labelenumi{\arabic{enumi}.}
\tightlist
\item
  The user starts in \textbf{Bulk OCR}, either by uploading files through the browser or by selecting server-side PDFs already copied into the working directory.
\item
  The OCR page sends the selected files to the OCR service and stores the returned artifacts as document folders containing \texttt{ocr\_result.json}, page markdown files, and image crops.
\item
  The user then moves to \textbf{LLM Processing}, which loads one OCR-expanded document at a time.
\item
  Inside \textbf{LLM Processing}, the user selects a page range rather than always processing the full report. This is important because the extraction workflow is designed around page batches.
\item
  The user chooses one provider family for the extraction stage: \textbf{OpenRouter}, \textbf{LM Studio / OpenAI-compatible}, or \textbf{Ollama}.
\item
  The selected page batch is sent to the chosen LLM provider, and the returned output is parsed into the ESG record schema.
\item
  Structured records are appended into \texttt{results/esg\_records.json}, which becomes the main evidence layer for later analysis.
\item
  The extracted text layer is then passed into the \textbf{ClimateBERT} T1 stage, which acts as a downstream comparison or benchmark layer rather than as the first ingestion step.
\item
  ClimateBERT-style predictions are saved into \texttt{results/predictions.json} or resumable T1 artifacts such as \texttt{results/t1\_results.jsonl}.
\end{enumerate}

\subsection{LaTeX Figure Version}\label{latex-figure-version}

The same operational workflow can also be represented in a LaTeX-ready figure block for the final thesis.

% \begin{figure}[H]
% \begin{center}
%   \safeimage[width=\textwidth]{appendix_workflow_bulkocr_llm_climatebert}
% \end{center}
\caption{Appendix workflow from Bulk OCR to page-range-based LLM extraction and downstream ClimateBERT processing. The implemented flow begins with OCR expansion, continues with provider-specific ESG extraction over selected page ranges, and ends with ClimateBERT comparison over the extracted text layer.}
\alt{Workflow diagram showing user upload or file selection in Bulk OCR, OCR storage into thesis dataset folders, page-range selection in LLM Processing, choice among OpenRouter, LM Studio, or Ollama, storage of structured ESG records, and downstream ClimateBERT T1 processing.}
\label{fig:appendix_bulkocr_llm_climatebert_workflow}
\end{figure}

\subsection{Interpretation}\label{interpretation}

This diagram shows that ClimateBERT processing is not the first-stage ingestion tool. Instead, it is a downstream comparison or benchmark layer applied after OCR expansion and after page-range-based ESG extraction. That distinction is methodologically important because the thesis does not compare ClimateBERT directly against raw PDFs. It compares ClimateBERT-style predictions against the extracted text layer produced by the broader ESG pipeline.

\subsection{Tone-Prompt Family Visualization and Prompt Reasoning}\label{tone-prompt-family-visualization-and-prompt-reasoning}

The repository contains six thesis-facing tone prompts under \texttt{prompt/}:

\begin{itemize}
\tightlist
\item
  \texttt{tone\_zero\_shot\_english.md}
\item
  \texttt{tone\_zero\_shot\_indonesian.md}
\item
  \texttt{tone\_few\_shot\_english.md}
\item
  \texttt{tone\_few\_shot\_indonesian.md}
\item
  \texttt{tone\_chain\_of\_thought\_english.md}
\item
  \texttt{tone\_chain\_of\_thought\_indonesian.md}
\end{itemize}

These prompt files share the same extraction target, but they differ in how much reasoning structure they impose before the final JSON output. The prompt family is summarized below so that the appendix makes clear that prompt variation is not cosmetic. It is one of the main experimental levers of the thesis.

\subsubsection{Prompt-Family Logic}\label{prompt-family-logic}

% \begin{figure}[H]
% \begin{center}
%   \safeimage[width=\textwidth]{appendix_tone_prompt_family_logic}
% \end{center}
% \caption{Tone-prompt family logic. This figure shows how the zero-shot, few-shot, and chain-of-thought prompt families share the same ESG input target but differ in reasoning scaffolds before producing structured JSON output.}
% \alt{Flowchart showing an ESG text input branching into zero-shot, few-shot, and chain-of-thought prompt families, each with distinct reasoning supports that converge into the same structured JSON output.}
% \label{fig:appendix_tone_prompt_family_logic}
% \end{figure}

\subsubsection{Prompt-Family Summary Table}\label{prompt-family-summary-table}

\begin{longtable}[]{@{}
  >{\raggedright\arraybackslash}p{(\columnwidth - 8\tabcolsep) * \real{0.2000}}
  >{\raggedright\arraybackslash}p{(\columnwidth - 8\tabcolsep) * \real{0.2000}}
  >{\raggedright\arraybackslash}p{(\columnwidth - 8\tabcolsep) * \real{0.2000}}
  >{\raggedright\arraybackslash}p{(\columnwidth - 8\tabcolsep) * \real{0.2000}}
  >{\raggedright\arraybackslash}p{(\columnwidth - 8\tabcolsep) * \real{0.2000}}@{}}
\caption{Tone-prompt family summary.}\label{tab:tone-prompt-family-summary}\\
\toprule\noalign{}
\begin{minipage}[b]{\linewidth}\raggedright
Prompt file
\end{minipage} & \begin{minipage}[b]{\linewidth}\raggedright
Prompt style
\end{minipage} & \begin{minipage}[b]{\linewidth}\raggedright
Main reasoning structure
\end{minipage} & \begin{minipage}[b]{\linewidth}\raggedright
Intended advantage
\end{minipage} & \begin{minipage}[b]{\linewidth}\raggedright
Main risk
\end{minipage} \\
\midrule\noalign{}
\endhead
\bottomrule\noalign{}
\endlastfoot
\texttt{tone\_zero\_shot\_english.md} & Zero-shot, English & Direct instructions, label list, tone priority rules, JSON schema & Strong baseline for testing whether explicit definitions alone are enough & May remain too shallow for ambiguous governance or mixed-tone text \\
\texttt{tone\_zero\_shot\_indonesian.md} & Zero-shot, Indonesian & Same logic as English version, but adapted to Indonesian cue words and examples of commitment/action/outcome distinction & Better alignment with Indonesian report language and bilingual segments & Longer instruction set can still be ignored when outputs drift \\
\texttt{tone\_few\_shot\_english.md} & Few-shot, English & Adds worked examples for commitment, outcome, and risk-oriented none tone & Demonstration-based anchoring for label and tone boundaries & Examples may overfit narrow patterns and suppress extraction breadth \\
\texttt{tone\_few\_shot\_indonesian.md} & Few-shot, Indonesian & Adds Indonesian worked examples and local wording for target, risk, and measured result & Better grounding for Indonesian narrative style and common lexical cues & Example dependence may still fail on governance-heavy or non-example structures \\
\texttt{tone\_chain\_of\_thought\_english.md} & Chain-of-thought style, English & Hidden internal steps for segmentation, tone assignment, sentiment assignment, and final schema fill & Forces the model to separate reasoning stages before output & Longer prompt can increase verbosity pressure or schema drift in weaker models \\
\texttt{tone\_chain\_of\_thought\_indonesian.md} & Chain-of-thought style, Indonesian & Most explicit reasoning scaffold, local cue words, tone priority logic, and concise reasoning guidance & Best fit for difficult Indonesian and mixed-language extraction conditions & Most complex prompt, so weaker backends may fail to follow every instruction consistently \\
\end{longtable}

\subsubsection{Why the Prompt Reasoning Matters}\label{why-the-prompt-reasoning-matters}

The central prompt-design idea is not just to ask for JSON, but to force the model to reason about disclosure function before assigning sentiment. In this thesis, that means the prompt must first distinguish:

\begin{enumerate}
\def\labelenumi{\arabic{enumi}.}
\tightlist
\item
  whether a segment is explicitly ESG-relevant;
\item
  what aspect and label family it belongs to;
\item
  whether the statement is a commitment, action, or outcome;
\item
  only after that, whether the statement carries positive, negative, neutral, or none sentiment.
\end{enumerate}

This ordering matters because many sustainability disclosures are future-oriented promises written in positive language. If the prompt does not explicitly disentangle tone from sentiment, the model can collapse a promise into a falsely positive achieved result. The tone prompts therefore use reasoning scaffolds to reduce this specific validity error.

\subsubsection{Prompt Reasoning Description by Family}\label{prompt-reasoning-description-by-family}

\textbf{Zero-shot reasoning.} The zero-shot prompts rely on explicit definitions, priority rules, and schema constraints without demonstration examples. Their reasoning philosophy is that tone extraction should remain possible from a strong instruction set alone. In methodological terms, zero-shot prompts act as the cleanest test of whether the schema itself is intelligible to the model.

\textbf{Few-shot reasoning.} The few-shot prompts add short worked examples so the model sees what commitment, outcome, and climate-risk cases should look like in JSON form. Their reasoning philosophy is analogical: the model is guided to map new inputs to demonstrated output patterns. This is useful for ambiguous ESG phrasing, but it can also narrow the extraction behavior if the examples are too stereotyped.

\textbf{Chain-of-thought-style reasoning.} The chain-of-thought prompts explicitly instruct the model to segment text, identify ESG relevance, assign tone, then assign sentiment, while keeping the actual reasoning hidden from the final response. Their reasoning philosophy is procedural: a difficult extraction problem is decomposed into smaller internal steps so that tone and sentiment are less likely to be conflated. This is the most thesis-aligned prompt family because the research question depends on stable tone disentanglement rather than generic JSON completion alone.

\subsubsection{Appendix Interpretation}\label{appendix-interpretation}

This prompt family should be read as a controlled prompt-engineering ladder rather than as six unrelated files. The zero-shot prompts test direct schema interpretability. The few-shot prompts test example anchoring. The chain-of-thought prompts test whether stepwise internal reasoning improves tone stability. That is why prompt variation in Chapter 4 is analytically meaningful: it probes whether the thesis contribution depends only on model choice or also on reasoning structure embedded in the prompt itself.

\section{JSON Data Examples}\label{json-data-examples}

The repository uses JSON extensively for OCR outputs, extraction records, dashboard artifacts, model caches, logs, and workflow decisions. This appendix section provides representative examples of the actual JSON structures used by the implemented system.

\subsection{OCR Output JSON}\label{ocr-output-json}

Path example:

\begin{itemize}
\tightlist
\item
  \texttt{data/thesis\_dataset/2017\ Sustainability\ Report\ PT\ Bank\ Permata\ Tbk\ (Rev\ May2018)\_0\_pdf/ocr\_result.json}
\end{itemize}

Purpose:

\begin{itemize}
\tightlist
\item
  stores page-level OCR output
\item
  preserves extracted markdown text
\item
  stores image coordinates and embedded image payloads
\item
  records OCR model and usage metadata
\end{itemize}

Example:

\begin{verbatim}
{
  "pages": [
    {
      "index": 0,
      "markdown": "PermataBank\n\nLaporan Keberlanjutan 2017...",
      "images": [
        {
          "id": "img-0.jpeg",
          "top_left_x": 48,
          "top_left_y": 260,
          "bottom_right_x": 667,
          "bottom_right_y": 777,
          "image_base64": "data:image/jpeg;base64,..."
        }
      ]
    }
  ],
  "model": "...",
  "document_annotation": "...",
  "usage_info": "..."
}
\end{verbatim}

\subsection{ESG Extraction Run JSON}\label{esg-extraction-run-json}

Path example:

\begin{itemize}
\tightlist
\item
  \texttt{results/esg\_records.json}
\end{itemize}

Purpose:

\begin{itemize}
\tightlist
\item
  stores T3 LLM extraction runs
\item
  preserves prompt, model, page target, records, and failure state
\item
  acts as the main structured ESG evidence store
\end{itemize}

Success example:

\begin{verbatim}
{
  "timestamp": "2026-06-03T18:38:40Z",
  "model": "nvidia/nemotron-3-nano-30b-a3b:free",
  "target": "Sustainability_Report_PT_Sentul_City_Tbk_2025_pdf/batch_7",
  "target_pages": "page_0012.md, page_0013.md",
  "prompt": "tone_chain_of_thought_indonesian.md",
  "ok": true,
  "records": [
    {
      "text": "Elevating Lives, Building Sustainable Growth",
      "aspect": "strategic vision",
      "labels": ["climate-d", "climate-commitment"],
      "esg": "E",
      "tone": "commitment",
      "sentiment": "positive",
      "sentiment_score": 1,
      "reasoning": "The phrase expresses a forward-looking ..."
    }
  ]
}
\end{verbatim}

Error example:

\begin{verbatim}
{
  "timestamp": "2026-06-03T18:37:59Z",
  "model": "nvidia/nemotron-nano-12b-v2-vl:free",
  "target": "Bank Neo Sustainable Report/batch_27",
  "target_pages": "page_0052.md, page_0053.md",
  "prompt": "tone_few_shot_indonesian.md",
  "ok": false,
  "records": [],
  "error": "OpenRouter returned HTTP 401: {\"error\":{\"message..",
  "error_type": "unknown",
  "raw_output": "",
  "background_job_id": "llm_processing_bg_20260603T183739Z_855dc4"
}
\end{verbatim}

\subsection{ClimateBERT Result JSON}\label{climatebert-result-json}

Path example:

\begin{itemize}
\tightlist
\item
  \texttt{results/climatebert\_results.json}
\end{itemize}

Purpose:

\begin{itemize}
\tightlist
\item
  stores ClimateBERT or remote-space comparison runs
\item
  keeps the original input text
\item
  preserves raw multi-model response text
\item
  supports later parsing into label-level comparison tables
\end{itemize}

Example:

\begin{verbatim}
{
  "timestamp": "2026-03-26T15:43:20.151168Z",
  "mode": "remote_space",
  "space_url": "https://darisdzakwanhoesien-climatebert..",
  "input_text": "Company's media and entertainment ..",
  "response_raw": "### ...\n### netzero-reduction\n•",
  "response_parsed": null
}
\end{verbatim}

\subsection{Other JSON Families Used by the Repository}\label{other-json-families-used-by-the-repository}

Additional JSON artifact families in the repository include the following. These are not all equally central to the thesis argument, but each one supports a specific part of the executable research workflow.

\begin{itemize}
\item
  \texttt{results/data/mapping.json} Purpose: stores label-normalization mappings used to standardize ESG and sentiment terminology. Example content includes mappings such as \texttt{"environmental":\ "E"}, \texttt{"governance":\ "G"}, and \texttt{"lingkungan":\ "E"}. This file acts as a normalization bridge between raw extracted labels and the thesis-facing schema.

\item
  \texttt{results/ground\_truth.json} Purpose: stores manually entered or review-oriented ground-truth style records. The sampled structure includes a timestamp, model name, source, input text, and a nested \texttt{result} object. In the current repository state, some entries still preserve model-side errors, which is useful because it shows that the reference-building layer keeps incomplete or failed cases visible rather than silently discarding them.
\item
  \texttt{results/predictions.json} Purpose: stores prediction outputs, especially from ClimateBERT-style or T1 comparison runs. The current file contains long source text segments together with model identifiers and prediction payloads. In practice, this file acts as a raw comparison layer before more compact summary tables are derived.

\item
  \texttt{results/t1\_results.json} Purpose: stores serialized T1-stage prediction results in JSON form. These records are similar to \texttt{predictions.json}, but they are part of the more formal T1 artifact family used by the benchmark side of the workflow. The sampled entries show model names, original text, and nested result objects.

\item
  \texttt{results/t2\_results.json} Purpose: stores T2-stage ABSA-style outputs. The sampled entries contain both a \texttt{rule\_based} block and a \texttt{hybrid} block. This is important because it shows that T2 is not one monolithic classifier. It compares rule-based aspect/polarity/tone output against a hybrid contextual model that also records ontology alignment and summary metrics.
\end{itemize}

Together, these JSON families support benchmark generation, ontology mapping, workflow documentation, summarization inputs, transfer-learning dataset reporting, and interactive tooling. They form an important part of the repository's reproducibility and audit layer even when they are not all directly cited in the main thesis chapters.
